% Generated draft: requires mathematical and source review.
\documentclass[11pt]{article}
\usepackage[margin=25mm]{geometry}
\usepackage{fontspec}
\setmainfont{Latin Modern Roman}
\usepackage{amsmath,amssymb,mathtools,unicode-math}
\setmathfont{Latin Modern Math}
\usepackage{enumitem,needspace,xurl,hyperref}
\setlength{\parindent}{0pt}
\setlength{\parskip}{5pt}
\setlength{\emergencystretch}{3em}
\newcommand{\N}{\mathbb{N}}
\newcommand{\Z}{\mathbb{Z}}
\newcommand{\Q}{\mathbb{Q}}
\newcommand{\R}{\mathbb{R}}
\newcommand{\C}{\mathbb{C}}
\newcommand{\F}{\mathbb{F}}
\newcommand{\A}{\mathbb{A}}
\newcommand{\PP}{\mathbb{P}}
\newcommand{\E}{\mathbb{E}}
\newcommand{\eps}{\varepsilon}
\newcommand{\dd}{\,\mathrm{d}}
\begin{document}
\section*{UnsolvedMath ID: 30003205}
\textbf{Source problem code:} catalogue entry
\subsection{Definitions and mathematical statement}
Let \(R\) denote the hyperfinite factor of type \(\mathrm{II}_1\) referred to in the supplied conjecture, and let \(\tau\) be its normalized faithful normal trace, so that \(\tau(1)=1\). A factor is a von Neumann algebra whose center consists of scalar multiples of its identity. Type \(\mathrm{II}_1\) means that the factor is finite and has no nonzero minimal projections. Hyperfiniteness means that the algebra is the weak operator closure of a directed union of finite-dimensional unital \(*\)-subalgebras.

Define
\[
\mathcal U(R)=\{u\in R:u^*u=uu^*=1\},\qquad
\|x\|_2=\tau(x^*x)^{1/2},\qquad
 d_2(u,v)=\|u-v\|_2.
\]
A subgroup \(\Gamma\leq\mathcal U(R)\) is uniformly discrete for this metric if there is a real number \(\delta>0\) such that
\[
\|g-h\|_2\geq\delta
\quad\text{for all distinct }g,h\in\Gamma.
\]
By unitary invariance of the trace norm, this is equivalent to \(\|g-1\|_2\geq\delta\) for every \(g\in\Gamma\setminus\{1\}\). A free nonabelian subgroup means a subgroup isomorphic to a free group on at least two generators.

\textbf{Conjecture (A. Thom, attributed as Andreas Thom in the supplied verification note).} The group \(\mathcal U(R)\) has no uniformly discrete free nonabelian subgroup with respect to \(d_2\).

Equivalently, let \(F_2=\langle a,b\rangle\) be the abstract free group on two generators, with identity \(e\). The assertion is
\[
\forall\text{ injective homomorphisms }\rho:F_2\longrightarrow\mathcal U(R),\quad
\forall\varepsilon>0,\quad
\exists w\in F_2\setminus\{e\}:\quad
0<\|\rho(w)-1\|_2<\varepsilon.
\]
The rank-two formulation suffices because every free nonabelian group contains a subgroup isomorphic to \(F_2\), and a subgroup of a uniformly discrete group is uniformly discrete.

\textit{Unresolved scope qualification.} The source uses the phrase ``the hyperfinite \(\mathrm{II}_1\) factor'' without explicitly stating separability. Its usual interpretation is the hyperfinite \(\mathrm{II}_1\) factor with separable predual. The supplied record does not specify whether any broader, nonseparable formulation is intended; no such extension is asserted here.
\subsection{Short English statement}
Thom's conjecture says that any algebraically free pair of unitaries in the specified hyperfinite \(\mathrm{II}_1\) factor generates nonidentity elements arbitrarily close to the identity in the normalized trace \(2\)-norm. The record labels the conjecture open; this is supplied curation, not an independently verified status assessment. The separability convention requires review.
\subsection{Sources}
\begin{itemize}
\item Local UnsolvedMath registry: supplied complete source record, permanent ID \(30003205\), exact source problem code \texttt{catalogue entry}, in the Oberwolfach Reports Open Problems collection. The original statement attributes the conjecture to A. Thom; the supplied verification note identifies him as Andreas Thom. The proposed year is recorded as 2016.
\item Supplied report citation, with its mathematical title typography normalized: \emph{\(C^*\)-Algebras (2017)}, Oberwolfach Reports, DOI \texttt{10.4171/owr/2016/40}, pp.\ 2269--2345. \url{https://doi.org/10.4171/owr/2016/40}. The supplied title year is 2017, whereas the DOI, proposed year, and source year tag contain 2016; the record does not explain this difference.
\item Supplied literature source: Carderi and Thom, \emph{Exotic Polish group limits}. \url{https://www.numdam.org/article/AIF_2018__68_1_257_0.pdf}.
\item Supplied literature source: \emph{Waterloo operator algebra open problems}. \url{https://www.math.uwaterloo.ca/~j37peter/problems.html}.
\item Supplied literature source: \emph{OWR Operator Algebras report}. \url{https://ems.press/content/serial-article-files/46645?nt=1}.
\item Supplied literature assessment dated exactly \texttt{2026-08-21}: status recorded as open. The assessment says that the cited theorem is conditional, imposes constraints on a possible counterexample or proof, and does not resolve the underlying discrete free subgroup conjecture. This is supplied dated curation derived from the OWR material, not an independently verified literature assessment. The background expressly calls for independent verification of time-sensitive status and cited claims.
\item The registry labels the statement \texttt{corrected\_verified} and reports that the official report was checked and the attributional heading removed from the mathematical statement. These are supplied provenance claims; no source visit or independent verification is claimed in this draft.
\end{itemize}
\end{document}
